% =============================================================================
% 学术海报模板
% FormalMath项目学术写作支持系统
% 适用：学术会议海报展示
% 尺寸：A0 (841mm × 1189mm)
% =============================================================================

\documentclass[a0paper,portrait]{baposter}

% ==================== 宏包 ====================
% 中文支持
\usepackage{ctex}

% 数学
\usepackage{amsmath,amssymb,amsthm}

% 图形
\usepackage{graphicx}
\usepackage{tikz}
\usepackage{tikz-cd}
\usetikzlibrary{arrows,shapes}

% 颜色
\usepackage{xcolor}

% 字体
\usepackage{relsize}

% 多栏
\usepackage{multicol}

% ==================== 颜色定义 ====================
\definecolor{maincolor}{RGB}{31, 78, 121}
\definecolor{accentcolor}{RGB}{192, 80, 77}
\definecolor{boxcolor}{RGB}{240, 240, 240}
\definecolor{headercolor}{RGB}{200, 220, 240}

% ==================== 海报设置 ====================
\begin{document}

\begin{poster}{
    % 海报选项
    grid=false,
    columns=2,
    colspacing=1em,
    headerheight=0.12\textheight,
    background=user,
    bgColorOne=white,
    headershape=rounded,
    headershade=plain,
    headerColorOne=maincolor,
    headerColorTwo=maincolor,
    headerfont=\Large\bfseries\color{white},
    boxshade=plain,
    boxColorOne=boxcolor,
    boxColorTwo=boxcolor,
    textborder=faded,
    boxheaderheight=2em,
    headerborder=open,
    boxpadding=1em
}
% ==================== 标题区 ====================
{
    % 标题
    \bfseries\Huge 海报标题：研究的简短描述\\[0.3em]
    \Large 作者姓名\textsuperscript{1}, 合作者\textsuperscript{2}\\[0.2em]
    \large \textsuperscript{1}学校/院系，\textsuperscript{2}合作机构\\[0.2em]
    \normalsize 联系邮箱：email@example.com
}
% ==================== 左侧列 ====================
{
    % 1. 引言
    \begin{posterbox}{1. 引言}
        \textbf{研究背景：}

        简要介绍研究背景，说明为什么这个问题重要。限制在100字以内。

        \vspace{0.5em}

        \textbf{核心问题：}

        明确陈述研究要解决的具体问题。可以使用项目符号：
        \begin{itemize}
            \item 问题1
            \item 问题2
        \end{itemize}
    \end{posterbox}

    % 2. 预备知识/方法
    \begin{posterbox}{2. 方法与工具}
        \textbf{主要工具：}

        介绍使用的主要数学工具或方法。

        \vspace{0.5em}

        \begin{center}
        \begin{tikzpicture}[scale=1.5]
            % 方法示意图
            \node[draw, rounded corners, fill=white] (A) at (0,0) {输入};
            \node[draw, rounded corners, fill=white] (B) at (3,0) {处理};
            \node[draw, rounded corners, fill=white] (C) at (6,0) {输出};
            \draw[->, thick] (A) -- (B);
            \draw[->, thick] (B) -- (C);
        \end{tikzpicture}
        \end{center}

        \vspace{0.5em}

        流程说明文字...
    \end{posterbox}

    % 3. 关键定义
    \begin{posterbox}{3. 关键定义}
        \textbf{定义 3.1} 关键概念的定义。

        \[
        f(x) = \sum_{n=0}^{\infty} a_n x^n
        \]

        \vspace{0.5em}

        \textbf{记号说明：}

        \begin{tabular}{ll}
            $\mathbb{R}$ & 实数集 \\
            $\mathbb{C}$ & 复数集 \\
            $GL_n$ & 一般线性群
        \end{tabular}
    \end{posterbox}
}
% ==================== 右侧列 ====================
{
    % 4. 主要结果
    \begin{posterbox}{4. 主要结果}
        \textbf{定理 4.1}（主要定理）

        陈述主要定理。使用框突出重要结果。

        \begin{center}
        \fbox{\parbox{0.9\linewidth}{
            \centering
            \textbf{核心等式}\\[0.5em]
            $\displaystyle E = mc^2$
        }}
        \end{center}

        \vspace{0.5em}

        \textbf{证明概要：}

        简要说明证明的主要步骤...
    \end{posterbox}

    % 5. 例子与应用
    \begin{posterbox}{5. 例子与应用}
        \textbf{示例 5.1}

        一个具体的例子说明理论的应用。

        \vspace{0.5em}

        \begin{center}
        \begin{tikzpicture}[scale=2]
            \draw[->] (-1.5,0) -- (1.5,0) node[right] {$x$};
            \draw[->] (0,-0.5) -- (0,1.5) node[above] {$y$};
            \draw[thick, blue, domain=-1:1, samples=50] plot (\x, {\x*\x});
            \node at (1.2,1) {$y=x^2$};
        \end{tikzpicture}
        \end{center}

        图1：示例函数图像
    \end{posterbox}

    % 6. 结论与参考文献
    \begin{posterbox}{6. 结论与参考文献}
        \textbf{主要结论：}
        \begin{itemize}
            \item 结论1
            \item 结论2
            \item 未来研究方向
        \end{itemize}

        \vspace{0.5em}

        \textbf{参考文献：}

        \begin{enumerate}[{[1]}]
            \item J. Smith, \textit{Important Results}, J. Math. 2020.
            \item A. Jones, \textit{Related Work}, Ann. Math. 2021.
            \item 基金资助信息（如有）
        \end{enumerate}

        \vspace{1em}

        \begin{center}
        \textbf{欢迎扫描二维码获取论文预印本！}

        \vspace{0.5em}

        % 放置二维码
        \framebox(2cm,2cm){二维码}
        \end{center}
    \end{posterbox}
}

\end{poster}
\end{document}
% =============================================================================
% 使用说明：
% 1. 使用XeLaTeX编译
% 2. 需要安装baposter宏包
% 3. 打印前检查分辨率（建议150-300 DPI）
% 4. 内容简洁，突出重点
% 5. 字体大小确保1米外可读
% 6. 准备1分钟和5分钟两种介绍版本
% =============================================================================
